\documentclass[11pt]{article}
\usepackage[T1]{fontenc}
\usepackage[utf8]{inputenc}
\usepackage[margin=1in]{geometry}
\usepackage{amsmath,amssymb}
\usepackage{booktabs}
\usepackage{microtype}
\usepackage[hidelinks]{hyperref}
\newcommand{\E}{\mathbb E}
\newcommand{\Pb}{\mathbb P}
\newcommand{\Ind}[1]{\mathbf 1\{#1\}}
\newcommand{\K}{K_\rho}
\newcommand{\xx}{\mathbf x}
\allowdisplaybreaks[1]
\title{Two-Seller Competition with Partial Information}
\author{Progress note --- transcription of handwritten notes}
\date{October 11, 2026}

\begin{document}
\maketitle

This note transcribes the model in the handwritten notes. We use the
coordinate order $(x_A,x_B)$ throughout and write seller-specific
expressions with a consistent index $j\in\{A,B\}$. Questions in the
notes about the boundary and the stationary distribution remain open.

\section{Environment, information, and timing}

Consider two sellers, indexed by $j\in\{A,B\}$, each with two products,
$i\in\{1,2\}$. Product $i$ generates a successful outcome with
probability $p_i$. The model incorporates forgetting with parameter
$\rho\in(0,1)$. The state at time $t$ is
\begin{equation}\label{eq:state}
\xx_t=(x_{A,t},x_{B,t}),
\end{equation}
where $x_{j,t}$ summarizes seller $j$'s successes under forgetting.
Each seller observes only its own outcomes and state $x_{j,t}$, not
the other seller's state. Write $-j$ for seller $j$'s opponent.

Following the restriction in the notes, we work with
\begin{equation}\label{eq:state-space}
\rho\in\left(\frac12,1\right),\qquad
\K=[0,1],\qquad \xx_t\in\K^2.
\end{equation}
This restriction is intended to avoid the Cantor-set issue associated
with smaller values of $\rho$. The notes ask whether the proposed
state representation amounts to working on the boundary of the original
forgetting model; this interpretation remains to be clarified.

\paragraph{Timing.}
Let $y_t\in\{A,B\}$ denote the customer's choice, $i_t\in\{1,2\}$ the
selected seller's product, and $u_t\in\{0,1\}$ the realized outcome.
Within each period:
\begin{enumerate}
\item The customer chooses a seller $j$.
\item The selected seller chooses a product $i$.
\item The outcome is realized, with success probability $p_i$, and
the selected seller's state is updated.
\end{enumerate}
In particular,
\begin{equation}\label{eq:outcome}
\Pb(u_t=1\mid y_t=j,i_t=i,\xx_t)=p_i.
\end{equation}

\paragraph{Customer choice.}
The selection probabilities are
\begin{equation}\label{eq:demand}
\Pb(y_t=j\mid\xx_t)=q_j(\xx_t)
=\Lambda(x_{j,t}-x_{-j,t}),
\end{equation}
where $\Lambda:[-1,1]\to[0,1]$ is continuous and increasing, with
\begin{equation}\label{eq:demand-symmetry}
\Lambda(-z)=1-\Lambda(z).
\end{equation}
Consequently, $q_A(\xx)+q_B(\xx)=1$.

\clearpage
\section{Strategies and the joint transition kernel}

\paragraph{State updates.}
For a selected seller with current state $x$, define
\begin{equation}\label{eq:scalar-transitions}
T^s(x)=\rho x+(1-\rho),\qquad T^f(x)=\rho x,
\end{equation}
where $s$ and $f$ denote success and failure. The unselected seller's
state remains unchanged.

\paragraph{Strategies.}
Fix a strategy profile $\sigma=(\sigma_A,\sigma_B)$, where
\begin{equation}\label{eq:strategies}
\sigma_j:\K\longrightarrow\Delta(\{1,2\}).
\end{equation}
Here $\Delta(\{1,2\})$ is the set of probability distributions over
the two products, and $\sigma_j(i\mid x_j)$ is the probability that
seller $j$ chooses product $i$ when selected at its own state $x_j$.
The resulting average success probability is
\begin{equation}\label{eq:average-success}
\bar p_j(x_j)=\sum_{i=1}^2\sigma_j(i\mid x_j)p_i.
\end{equation}

From a joint state $\xx=(x_A,x_B)$, the four possible events have the
following next states and probabilities:
\begin{center}
\renewcommand{\arraystretch}{1.3}
\begin{tabular}{lll}
\toprule
Event & Next state & Probability \\
\midrule
$A$ selected, success & $(T^s(x_A),x_B)$
  & $q_A(\xx)\bar p_A(x_A)$ \\
$A$ selected, failure & $(T^f(x_A),x_B)$
  & $q_A(\xx)(1-\bar p_A(x_A))$ \\
$B$ selected, success & $(x_A,T^s(x_B))$
  & $(1-q_A(\xx))\bar p_B(x_B)$ \\
$B$ selected, failure & $(x_A,T^f(x_B))$
  & $(1-q_A(\xx))(1-\bar p_B(x_B))$ \\
\bottomrule
\end{tabular}
\end{center}

Let $\delta_z$ denote the Dirac probability measure at $z$; for a
measurable set $E$, $\delta_z(E)=\Ind{z\in E}$. The joint transition
kernel, the continuous-state counterpart of a transition matrix, is
\begin{equation}\label{eq:joint-kernel}
\begin{aligned}
P^\sigma(\xx,\cdot)
={}&q_A(\xx)\Bigl[
  \bar p_A(x_A)\delta_{(T^s(x_A),x_B)}
 +(1-\bar p_A(x_A))\delta_{(T^f(x_A),x_B)}\Bigr]\\
&+(1-q_A(\xx))\Bigl[
  \bar p_B(x_B)\delta_{(x_A,T^s(x_B))}
 +(1-\bar p_B(x_B))\delta_{(x_A,T^f(x_B))}\Bigr].
\end{aligned}
\end{equation}

\section{Stationary distribution and expected demand}

Let $\pi^\sigma$ be an invariant probability distribution for
$P^\sigma$, satisfying
\begin{equation}\label{eq:stationary-distribution}
\pi^\sigma(E)=\int_{\K^2}P^\sigma(\xx,E)\,\pi^\sigma(d\xx)
\qquad\text{for every measurable }E\subseteq\K^2.
\end{equation}
As in the notes, we assume that this distribution exists and is unique.
The conditions that justify existence and uniqueness remain an open
question.

For seller $j$, define stationary expected demand conditional on its
own state by
\begin{equation}\label{eq:conditional-demand}
\begin{aligned}
\eta_j^\sigma(x)
&:=\E^{\pi^\sigma}\!\left[q_j(\xx_t)\mid x_{j,t}=x\right]\\
&=\int_{\K}\Lambda(x-x')\,
  \pi^\sigma(dx'\mid x_j=x).
\end{aligned}
\end{equation}
Here $\pi^\sigma(dx'\mid x_j=x)$ denotes the conditional distribution
of the opponent's state $x_{-j}=x'$ under $\pi^\sigma$.
Thus $\eta_j^\sigma(x)$ depends only on seller $j$'s own state: the
opponent's unobserved state is averaged out using its conditional
stationary distribution.

\section{Transition law conditional on a seller's own state}

Under this stationary demand, the one-period conditional distribution
of seller $j$'s next state is
\begin{equation}\label{eq:own-state-kernel}
\begin{aligned}
&\Pb^{\pi^\sigma}(x_{j,t+1}\in\cdot\mid x_{j,t}=x)\\
&\qquad=\eta_j^\sigma(x)\Bigl[
  \bar p_j(x)\delta_{T^s(x)}
 +(1-\bar p_j(x))\delta_{T^f(x)}\Bigr]
 +(1-\eta_j^\sigma(x))\delta_x.
\end{aligned}
\end{equation}
With probability $\eta_j^\sigma(x)$, seller $j$ is selected and its
state changes according to the outcome. Otherwise, its state remains
at $x$. Every term in this expression depends on the seller's own
state and the fixed strategy profile.

\paragraph{Technical clarification (not in the handwritten notes).}
Equation~\eqref{eq:own-state-kernel} is a one-period conditional law
under the stationary joint distribution. By itself, it does not
establish that a seller's own-state process is Markov with respect to
its entire observed history. Using it as a complete dynamic model
based only on $x$ would require a separate justification.

\end{document}
